
\chapter{Ie\c{s}an--Saint Venant Solution}\label{sec:sv}

The objective of this section is to briefly recall the family of general solutions to Saint-Venant's continuum problems \citep{iesan1987saintvenants,synge1945problem}. 
This is used in \cref{sec:timo} to identify a dimensionally reduced beam model that exactly represents these solutions. 


The exposition follows the approach of \cite{iesan1986saintvenants} and \cite{iesan1976saintvenantsb}, whereby the Saint Venant class is decomposed into two problems: that of torsion-bending-elongation (\cref{sec:iesan-pi}), and that of flexure (\cref{sec:iesan-pii}).
The developments are specialized in \cref{sec:iesan-poisson} to heterogeneous bodies with homogeneous Poisson ratio.

\begin{Quote}

Salient features of Iesan's treatment include:
\begin{itemize}
    \item The relaxed nature of Saint-Venant's class, including failure to exactly specify full boundary conditions in the continuum BVP, thus motivating the need for an additional postulate.
    \item Ie\c{s}an's postulates that uniquely specify Saint Venant's solution.
\end{itemize}
\end{Quote}

\section{Overview}
Infinitesimal motions of the body described in \cref{sec:std-continuum} are investigated, and attention is restricted to those composed heterogeneously of an isotropic material. 

The \emph{resultants} are the integrals of the traction \(-\StressTensor\,\FrameBasis\) on the root section \(\Section_0\):
\begin{equation}
\mathscr{F}(\hat{\bm u})=\int_{\Section_0}\StressTensor(\hat{\bm u})\,(-\FrameBasis)\dA,\qquad
\mathscr{M}(\hat{\bm u})=\int_{\Section_0}\bm{r}\times\big(\StressTensor(\hat{\bm u})\,(-\FrameBasis)\big)\dA.
\label{eq:def-resultants}
\end{equation}
where \(\StressTensor\;\) is the stress tensor. 
Saint Venant's solutions arise from  \emph{relaxed} elastostatic problems, where end tractions are not prescribed pointwise, but rather specified through these resultants. 
Neglecting to specify complete pointwise boundary data in this manner produces an ill-posed boundary-value problem which consequently lacks a unique solution \(\hat{\bm u}\):
\begin{equation}
\left\{
\begin{aligned}
\Div\,\StressTensor&=\mathbf{0} &&\text{in }\Body \\
\StressTensor\,\bm n&=\mathbf{0} && \text{on }\BodyLateralSurface\\
\mathscr{M}(\hat{\bm u}) &= \MomentRxn &&\text{on } \Section_0\\
\mathscr{F}(\hat{\bm u}) &= \bar{\bm{N}} & \\
\end{aligned}\right.
\qquad\text{with}\qquad 
%
\begin{aligned}
\StressTensor& \coloneq \lambda\,(\operatorname{tr}\StrainTensor)\,\mathbf{1}+2G\,\StrainTensor,\\
\StrainTensor& \coloneq \tfrac12\big(\nabla\hat{\bm u}+\nabla\hat{\bm u}^{\!\top}\big),
\end{aligned}
\label{eq:cauchy-space}
\end{equation}
In the treatment by \cite{iesan1986saintvenants} %resolves this indeterminacy in a systematic fashion that is well-suited to the present investigation. 
the family of Saint~Venant problems is decomposed into two classes, the extension-bending-torsion problem \(\mathrm{P}_{\rm I}(\AxialRxn, \MomentRxn)\) for prescribed axial force \(\AxialRxn\) and moment \(\MomentRxn\) resultants, and the flexure problem \(\mathrm{P}_{\mathrm{I\!I}}(\ShearRxn)\) for a prescribed transverse resultant \(\ShearRxn = \bar{F}_{\alpha} \,\FrameBasis_{\alpha}\):
\begin{equation}
\mathrm{P}_{\rm I}(\AxialRxn,\MomentRxn)\left\{
\begin{aligned}
\Div\,\StressTensor&=\mathbf{0} &&\text{in }\Body \\
\StressTensor\,\bm n&=\mathbf{0} &&\text{on }\BodyLateralSurface \\
\mathscr{M}(\hat{\bm u}) &= \MomentRxn &&\text{on } \Section_0 \\
\mathscr{F}(\hat{\bm u}) &= \AxialRxn\,\FrameBasis & \\
\end{aligned}\right.
%
\qquad\text{and}\qquad 
%
\mathrm{P}_{\mathrm{I\!I}}(\ShearRxn) \left\{
\begin{aligned}
\Div\,\StressTensor&=\mathbf{0} &&\text{in }\Body\\
\StressTensor\,\bm n&=\mathbf{0} &&\text{on }\BodyLateralSurface\\
\mathscr{M}(\hat{\bm u}) &= \mathbf{0} &&\text{on } \Section_0\\
\mathscr{F}(\hat{\bm u}) &= \bar{F}_{\alpha}\FrameBasis_{\alpha} & \\
\end{aligned}\right.
\label{eq:partial-elastostatic-equilibrium}
\end{equation}
% The problems \(\mathrm{P}_{\rm I}\) and \(\mathrm{P}_{\mathrm{I\!I}}\), introduced by \cite{iesan1986saintvenants}, decompose the Saint Venant class of problems into two uncoupled classes with a fundamental connection. 
\cite{iesan1986saintvenants} observes that mapping \(\hat{\bm u}\mapsto \hat{\bm u}'\) preserves lateral traction–free fields and transforms resultants as  \citep[Thm~2.1]{iesan1986saintvenants}\footnote{
The units for these identities appear paradoxical, and may be understood as the follows. 
Since \(\StressTensor(\hat{\bm u}')=(\StressTensor(\hat{\bm u}))'\), the resultants computed from \(\hat{\bm u}'\) carry an extra factor \(1/\text{length}\) relative to those of \(\hat{\bm u}\). 
Hence
\[
[\mathscr{F}(\hat{\bm u}')]=\text{force}/\text{length},\qquad
[\mathscr{M}(\hat{\bm u}')]=\text{force},
\]
so \(\mathscr{M}(\hat{\bm u}')\) indeed has the same units as \(\FrameBasis\times\mathscr{F}(\hat{\bm u})\) (force). Interpreting \(\mathscr{M}(\hat{\bm u}')\) as the \(\reflenx\)-derivative of the bending/torsion resultant, the identity is the continuum analogue of the 1D balance \(\mathscr{M}'=\FrameBasis\times \mathscr{F}\).
}
\[
\mathscr{F}(\hat{\bm u}')=\mathbf{0}, 
\qquad 
\mathscr{M}(\hat{\bm u}')=\mathscr{F}(\hat{\bm u})\times\FrameBasis,
\qquad
\FrameBasis\cdot\mathscr{M}(\hat{\bm u}')=0
\] 
Consequently, if a displacement \(\hat{\bm u}_{\mathrm{I\!I}}\) is a solution of the flexure problem \(\mathrm{P}_{\mathrm{I\!I}}(\ShearRxn)\), then its axial derivative \(\hat{\bm u}'_{\mathrm{I\!I}}\) is a solution of the extension-bending-torsion problem \(\mathrm{P}_{\rm I}(0,\,\FrameBasis\times\ShearRxn)\) \citep{iesan1987saintvenants}. % (Corollary~2.2). 
This connection motivates the program of Ie\c{s}an, whereby a particular solution to \(\mathrm{P}_{\rm I}\) is identified by prescribing the shape of \(\hat{\bm u}'\), and a solution to \(\mathrm{P}_{\mathrm{I\!I}}\), identified by prescribing the shape of \(\hat{\bm u}''\). 
These prescriptions select particular solutions within the relaxed class that are unique up to a rigid body field field, while leaving end tractions unspecified pointwise.

% \begin{remark}
% No restriction is imposed on the in–plane stress \((\FrameBasis_\alpha\otimes\FrameBasis_\beta)\,\StressTensor\) at this stage. 
% Components of interest are singled out as
% \[
% \AxialStress\coloneq\FrameBasis\cdot\StressTensor\,\FrameBasis\in\mathbb{R},
% \qquad
% \PointShearStress \coloneq(\FrameBasis_\beta\otimes\FrameBasis_\beta)\,\StressTensor\,\FrameBasis,
% \]
% % Ie\c{s}an’s construction recovers the Saint–Venant form \eqref{eq:saint-venant-stress}
% % \begin{equation*}
% % \StressTensor = \AxialStress\,\FrameBasis\otimes\FrameBasis
% % +\FrameBasis\otimes\bm{\tau}+\bm{\tau}\otimes\FrameBasis,
% % % \label{eq:saint-venant-stress}
% % \end{equation*}
% % i.e., the in–plane block vanishes as a consequence of the solution.
% \end{remark}

\section{Plane Warping}\label{sec:iesan-plane}
% \eqref{eq:def-eight-shapes}

\begin{subequations}
\begin{empheq}[left=\Pi(\FrameBasis_k)\empheqlbrace]{align}
\operatorname{div}\tilde{\StressTensor}^{(k)}+\nabla \big(\lambda\,(\FrameBasis_{k}\cdot(\bm{r} + \FrameBasis))\big)&=\mathbf{0},
&\text{in }\Section
\label{eq:ie-3-2-10} 
\\
%
\tilde{\StressTensor}^{(k)}\bm{\nu}&=-\FrameBasis_k\cdot(\bm{r}+\FrameBasis)\,\lambda\bm{\nu},
&\text{on }\SectionBoundary 
\label{eq:ie-3-2-11}
\end{empheq}
\label{eq:bvp-plane-shapes}
\end{subequations}

% \begin{subequations}
% \begin{align}
% \operatorname{div}\tilde{\StressTensor}^{(k)}+\nabla \big(\lambda\,(\FrameBasis_{k}\cdot(\bm{r} + \FrameBasis))\big)&=\mathbf{0},
% \qquad\text{in }\Section
% \label{eq:ie-3-2-10} 
% \\
% \tilde{\StressTensor}^{(k)}\bm{\nu}&=-\FrameBasis_k\cdot(\bm{r}+\FrameBasis)\,\lambda\bm{\nu},
% \qquad\text{on }\SectionBoundary 
% \label{eq:ie-3-2-11}
% \end{align}
% % \label{eq:bvp-plane-shapes}
% \end{subequations}
Define
\begin{equation}
\PlaneShape \coloneq \hat{\bm u}_{\Section}^{(\beta)}\otimes\FrameBasis_{\beta},
\qquad
\bar{\pi} \coloneq\operatorname{tr}\tilde{\StrainTensor}^{(3)},
\qquad
\PlaneShapeStrain
\coloneq \FrameBasis_\beta \operatorname{tr}\tilde{\StrainTensor}^{(\beta)},
\qquad \FrameBasis_\beta \in\Section,
\label{eq:def-Theta}
\end{equation}

\section{Problem I: extension--bending--torsion}\label{sec:iesan-pi}

\subsubsection*{Synopsis}
Problem \(\mathrm{P}_{\rm I}\) covers extension, bending, and torsion under resultants with \(F_\alpha=0\). The identifying kinematic postulate is that the first axial derivative is a rigid body motion,
\[
\hat{\bm u}'_{\rm I}= a_{\varepsilon}\,\FrameBasis
+\big(a_{\vartheta}\,\FrameBasis-\FrameBasis\times\bm{a}_{\Section}\big)\times\hat{\bm{x}},
\qquad
\bm{a}_{\Section}\in\mathbb{R}^2,\; a_{\varepsilon},a_{\vartheta}\in\mathbb{R},
\]
which fixes the \(\reflenx\)-dependence of \(\hat{\bm u}\) and reduces the equilibrium equations to planar problems on \(\Section\).

\subsubsection*{Statement}
Given a resultant axial force \(\AxialRxn\) and couple \(\MomentRxn\coloneq\TwistRxn\FrameBasis + \bar{M}_{\alpha}\,\FrameBasis_{\alpha}\) on \(\Section_{x_0}\),
determine \((\hat{\bm u},\StressTensor)\) such that:
\begin{equation}
\left\{
\begin{aligned}
\hat{\bm u}_{\rm I} &\in \mathrm{P}_{\rm I}(\AxialRxn,\,\MomentRxn) \\
\hat{\bm u}'_{\rm I}&= a_{\varepsilon}\,\FrameBasis
+\big(a_{\vartheta}\,\FrameBasis-\FrameBasis\times\bm{a}_{\Section}\big)\times\bm{x}.
\end{aligned}
\right.
\label{eq:iesan-problem-1}
\end{equation}

\subsubsection*{Solution}
Integrating \eqref{eq:iesan-problem-1} furnishes the particular solution:
\begin{equation}
\begin{aligned}
% \hat{\bm u}\{\bm a\} &\in \mathrm{P}_{\rm I}(\AxialRxn,\MomentRxn) \\
\hat{\bm u}_{\rm I}\{\bm a\}
&=\Big(-\tfrac12 \reflenx^2\,\mathbf{1}
+\reflenx\,\FrameBasis\otimes\bm{r}
+\PlaneShape
% -\nu\,\operatorname{dev}\bm{r}\otimes\bm{r}
\Big)\bm{a}_{\Section}
+(\reflenx\,\FrameBasis-\nu\,\bm{r})\,a_{\varepsilon}
+\big(\reflenx\,\FrameBasis\times\bm{r}+\WarpTwistIesan\,\FrameBasis\big)a_\vartheta ,
\end{aligned}
\label{eq:iesan-displacement-1}
\end{equation}
where the torsion warping function \(\WarpTwistIesan:\Section\rightarrow\mathbb{R}\) is determined by the Neumann problem:
\begin{equation}
\left\{
\begin{array}{rll}
\operatorname{div}_G\left(\left.\nabla \WarpSVP+\FrameBasis\times\PlanePosition\right.\right) &=0 & \text { in } \Section, \\
\nabla_{\bm{\nu}} \WarpSVP &=-\FrameBasis\times\PlanePosition \cdot \bm{\nu} & \text { on } \SectionBoundary, \\
\int_{\Section} \WarpSVP \volE &= 0
\end{array}
\right.
% \left\{
% \begin{aligned}
% \Delta \WarpTwistIesan&=0 && \text{in }\Section,\\
% \partial_{\bm{\nu}} \WarpTwistIesan&=-\,(\FrameBasis\times\PlanePosition)\cdot \bm{\nu} && \text{on } \SectionBoundary,\\
% \int_{\Section} \WarpTwistIesan\, \volE &=0, &&
% \end{aligned}
% \right.
\label{eq:bvp-warp-twist-iesan}
\end{equation}
\begin{Details}
\begin{equation*}
\left\{
\begin{aligned}
\operatorname{div}_G\!\big(\nabla \WarpTwistIesan+\FrameBasis\times\PlanePosition\big)&=0 && \text{in }\Section,\\
\partial_{\bm{\nu}} \WarpTwistIesan&=-\,(\FrameBasis\times\PlanePosition)\cdot \bm{\nu} && \text{on } \SectionBoundary,\\
\int_{\Section} \WarpTwistIesan\, \volE&=0, &&
\end{aligned}
\right.
% \label{eq:bvp-warp-twist-iesan}
\end{equation*}
\end{Details}
The linearized strain tensor under \eqref{eq:iesan-displacement-1} is:
\begin{equation}
\StrainTensor^{(a)}=
\big(a_{\varepsilon}+\bm{r}\cdot\bm{a}_{\Section}\big)\,(\FrameBasis\otimes\FrameBasis-\nu\,\oneS)
+\tfrac12\big(\FrameBasis\otimes\bm\varepsilon_\vartheta+\bm\varepsilon_\vartheta\otimes\FrameBasis\big),
\quad
\bm\varepsilon_\vartheta \coloneq a_\vartheta\big(\FrameBasis\times\bm{r}+\nabla\WarpTwistIesan\big),
\label{eq:iesan-strain-tensor-1}
\end{equation}
% and Saint–Venant components
% \begin{equation*}
% \bm{\tau}=G\big(\nabla\WarpTwistIesan+ \FrameBasis\times\bm{r}\big)a_\vartheta,
% \qquad
% \AxialStress=E\big(\bm{r}\cdot\bm{a}_{\Section}+a_{\varepsilon}\big).
% % \label{eq:iesan-stress-1}
% \end{equation*}
The integration constants \((\bm{a}_{\Section},a_{\varepsilon},a_\vartheta)\) are fixed by:
%
\begin{equation}
\begin{aligned}
\bm{a}_{\Section}\cdot \AxialCoupling+\AxialRigidity\, a_{\varepsilon} &=-\AxialRxn,\\
\CoupleRigidity\,\bm{a}_{\Section}+a_{\varepsilon}\,\AxialCoupling &=\MomentRxn\times\FrameBasis,\\
\TwistRigidity\,a_\vartheta&=-\TwistRxn,
\end{aligned}
\label{eq:2-15}
\end{equation}
where the axial rigidity \(\AxialRigidity\) and first moment vector \(\AxialCoupling\) are defined by:
\begin{equation}
\AxialRigidity \coloneq \int_\Section \Big[\lambda+2G+\lambda\bar{\pi}\Big]\dA
%\int \volE 
\qquad\text{and}\qquad
\AxialCoupling \coloneq \int_\Section \Big[(\lambda+2G)\,\bm{r}+\lambda\,\PlaneShapeStrain\Big]\dA
%\int \bm{r}\volE 
\label{eq:def-axial-rigidity}
\end{equation}
and the torsional rigidity \(\TwistRigidity\) and plane moment rigidity \(\CoupleRigidity\) are defined by:
\begin{equation}
\CoupleRigidity \coloneq \int_{\Section}\bm{r}\otimes\Big[(\lambda+2G)\,\bm{r}+\lambda\,\PlaneShapeStrain\Big]\dA
% \coloneq \int \bm{r}\otimes\bm{r}\volE
\qquad\text{and}\qquad
\TwistRigidity \coloneq \int_{\Section} \Big(r^2+(\FrameBasis\times\bm{r})\cdot\nabla\WarpTwistIesan\Big)\volG .
\label{eq:def-j-origin}
\end{equation}
Note that these are defined with respect to the unspecified coordinate origin. 

\section{Problem II: flexure}\label{sec:iesan-pii}

\subsubsection*{Synopsis}
Problem \(\mathrm{P}_{\mathrm{I\!I}}\) models flexure under a transverse resultant force \(\ShearRxn\) with \(\AxialRxn=0\) and \(\MomentRxn=\mathbf{0}\). 
The identifying kinematic postulate requires that the second axial derivative is a rigid motion,
\[
\hat{\bm u}_{\mathrm{I\!I}}''= b_{\varepsilon}\,\FrameBasis
+\big(b_{\vartheta}\,\FrameBasis-\FrameBasis\times\bm{b}_{\Section}\big)\times\bm{x},
\qquad
\bm{b}_{\Section}\in\mathbb{R}^2,\; b_{\varepsilon},b_{\vartheta}\in\mathbb{R}.
\]
It follows that \(\hat{\bm u}'\) is a \(\mathrm{P}_{\rm I}\) Saint–Venant field with bending moments \(\FrameBasis\times\ShearRxn\) %(Corollary~2.2)
, and hence
\(
\hat{\bm u}_{\mathrm{I\!I}}\{\bm b,\bm c\}=\int_0^{\reflenx}{\bm u}_{\rm I}\{\bm b\}\,\mathrm {~d}\reflenx+\bm{u}_{\rm I}\{\bm c\}+\tilde{\bm w}
\).


\subsubsection*{Statement}
Given a transverse shear resultant \(\ShearRxn\) on \(\Section_0\), determine \((\hat{\bm u}\{\bm{b},\bm{c}\},\,\StressTensor)\) such that:
%
\begin{subequations}
\begin{empheq}[left=\bar{\mathrm{P}}_{\mathrm{I\!I}}\empheqlbrace]{align}
\hat{\bm u}_{\mathrm{I\!I}}&\in\mathrm{P}_{\mathrm{I\!I}}(\ShearRxn) \\
\hat{\bm u}''_{\mathrm{I\!I}} &= b_{\varepsilon}\,\FrameBasis
+\big(b_{\vartheta}\,\FrameBasis-\FrameBasis\times\bm{b}_{\Section}\big)\times\bm{r},
\label{eq:iesan-postulate-2}
\end{empheq}
\label{eq:iesan-problem-2}
\end{subequations}


\subsubsection*{Solution}
Integrating \eqref{eq:iesan-problem-2} furnishes the particular solution:
\begin{equation}
\begin{aligned}
\oneS\hat{\bm u}_{\mathrm{I\!I}}(\reflenx,\bm{r})
&=-\tfrac16\,\reflenx^{3}\,\bm{b}_{\Section}
+\reflenx\,\nu\Big(\tfrac12 r^{2}\,\bm{b}_{\Section}-(\bm{b}_{\Section}\cdot\bm{r})\,\bm{r}\Big)
+\reflenx\,\nu\,(\CentroidLocation\cdot\bm{b}_{\Section})\,\bm{r}
+c_\vartheta\,\reflenx\,(\FrameBasis\times\bm{r}),\\
\FrameBasis\cdot\hat{\bm u}_{\mathrm{I\!I}}(\reflenx,\bm{r})
&=\tfrac12\,\reflenx^{2}\, (\bm{r}-\CentroidLocation)\cdot\bm{b}_{\Section}
+c_\vartheta\,\WarpTwistIesan(\bm{r})+\WarpShearShape{}(\bm{r}) \cdot \bm{b}_{\Section},
\end{aligned}
\label{eq:iesan-displacement-2}
\end{equation}
where 
\begin{equation}
\PlaneShearShape(\PlanePosition)\coloneq \PlaneShape-\PlaneAxialShape\otimes\CentroidLocation
\qquad\text{and}\qquad
\tilde{\PlaneShapeStrain}
\coloneq \PlaneShapeStrain - \bar{\pi}\,\CentroidLocation
\label{eq:def-plane-shear-shape}
\end{equation}
with
\begin{equation}
\CentroidLocation \coloneq \frac{1}{\AxialRigidity} \AxialCoupling
% \frac{1}{\AxialRigidity}\int \bm{r}\volE 
\label{eq:def-centroid}
\end{equation}
and the flexural warping shape \(\WarpShearShape{}:\Section\rightarrow\mathbb{R}^{2}\) is determined by the  Poisson–Neumann problem:
\begin{equation}
\left\{
\begin{aligned}
\Div \big(G\nabla(\WarpShearShape \cdot\bm{b}_{\Section})\big)&=\bm{f}_{\mathrm{I\!I}}\cdot\bm{b}_{\Section}
%-2\big(\bm{r}-\CentroidLocation\big)\cdot\bm{b}_{\Section} 
&& \text{in }\Section,\\
G\nabla_{\bm\nu}(\WarpShearShape\cdot\bm{b}_{\Section})&=-\bm{\nu}\cdot\PlaneShearShape\bm{b}_{\Section}
%\nu\,\bm\nu\cdot\Big(\operatorname{dev}\bm{r}\otimes\bm{r}-\bm{r}\otimes\CentroidLocation\Big)\bm{b}_{\Section}
&& \text{on } \SectionBoundary,\\
\int_\Section \WarpShearShape\cdot\bm{b}_{\Section}\,\volE&=0, &&
\end{aligned}
\right.
\label{eq:bvp-warp-shear-iesan}
\end{equation}
with:
\[
\bm{f}_{\mathrm{I\!I}}(\bm{r})
\coloneq -(\lambda+2G)\big(\bm{r}-\CentroidLocation\big)
-\operatorname{div}\big(G\,\PlaneShearShape\big)
-\lambda\,\tilde{\PlaneShapeStrain}\\
\]
% This problem has been considered by \cite{romano2012torsion} with the symbol \(\phi^{\rm SH}\), \cite{pilkey2003analysis} with the symbols \(\Phi,\Psi\), 
% \cite{lacarbonara2007solution} with the symbol \(\omega_i\),
% and \cite{hjelmstad1983seismic} with the symbol \(\Psi\). 

The strain under \eqref{eq:iesan-displacement-2} is
\begin{equation}
\StrainTensor(\hat{\bm u})
=(\PointAxialStrain\,\FrameBasis+\PointShearStrain)\symbun\FrameBasis
+\PlaneStrainTensor
\qquad\text{with}\qquad
% -\nu\,\varepsilon\,\oneS,
\begin{aligned}
\PointAxialStrain&=\reflenx\,(\bm{r}-\CentroidLocation)\cdot\bm{b}_{\Section},\\
\PointShearStrain
&= c_{\vartheta}\big(\nabla\WarpTwistIesan + \FrameBasis\times\bm{r}\big)
+ \PlaneShearShape\bm{b}_{\Section}
% -\nu\Big(\operatorname{dev}(\bm{r}\otimes\bm{r})-\bm{r}\otimes\CentroidLocation\Big)\bm{b}_{\Section}
+\nabla\big(\WarpShearShape{}\cdot\bm{b}_{\Section}\big),\\
\PlaneStrainTensor &= ...
\end{aligned}
\label{eq:iesan-strain-2}
\end{equation}
and the stress is
\begin{equation}
\begin{aligned}
\AxialStress(\hat{\bm u})&=E\,\reflenx\,(\bm{r}-\CentroidLocation)\cdot\bm{b}_{\Section},\\
\bm{\tau}(\hat{\bm u})
&=G\Big[(a_\vartheta+c_\vartheta)\big(\nabla\WarpTwistIesan+\FrameBasis\times\bm{r}\big)
+ \PlaneShearShape\bm{b}_{\Section}
% -\nu\Big(\operatorname{dev}(\bm{r}\otimes\bm{r})-\bm{r}\otimes\CentroidLocation\Big)\bm{b}_{\Section}
+\nabla\big(\WarpShearShape{}\cdot\bm{b}_{\Section}\big)\Big].
\end{aligned}
\label{eq:iesan-stress-2}
\end{equation}

% Differentiation in \(\reflenx\) links flexure to bending with \(\MomentRxn+\FrameBasis\times\ShearRxn=\mathbf{0}\).
The free parameters are determined from
\begin{equation}
\begin{aligned}
\CoupleRigidity \,\bm{b}_{\Section}+ \AxialCoupling\,b_\varepsilon &=-\ShearRxn,\\
\bm{b}_{\Section}\cdot\CentroidLocation+b_\varepsilon&=0,\\
b_\vartheta&=0,
\end{aligned}
\label{eq:2-18}
\end{equation}
and \(c_\vartheta\) is determined by the relaxed torsion boundary condition \(\FrameBasis\cdot\int \bm{r}\times\bm{\tau}\dA = 0\):
\begin{equation}
\,c_\vartheta
= \CenterShearLocation\cdot\bm{b}_{\Section}
% -\frac{1}{\Jsv}\FrameBasis\cdot\int_{\Section} \bm{r}\times \left(\nabla \WarpShearShape{}^{\top}
% + \PlaneShearShape \right)\bm{b}_{\Section} \dA 
% =-\frac{1}{\Jsv}\int_\Section\Big[(\FrameBasis\times\bm{r})\cdot\nabla\big(\WarpShearShape{}\cdot\bm{b}_{\Section}\big)
% +\tfrac12\,\nu\,r^2\,(\FrameBasis\times\bm{r})\cdot\bm{b}_{\Section}\Big]\dA .
\label{eq:iesan-2-22}
\end{equation}
where the \emph{shear point} \(\CenterShearLocation\in\Section\) is defined by:
\begin{equation}
\CenterShearLocation \coloneq \frac{1}{\TwistRigidity} \int_\Section G(\nabla\WarpShearShape{} + \PlaneShearShape^\top)\,\bm{r} \times \FrameBasis \dA.
\label{eq:def-shear-center}
\end{equation}
and \(\TwistRigidity\) is the torsion constant defined by \eqref{eq:def-j-origin}.

\begin{Details}
\[
\begin{aligned}
c_{\vartheta} &=- \frac{1}{\Jsv}\int_{\Section} \Big[(\FrameBasis\times \bm{r})\cdot \nabla \WarpShearIesan
+ \tfrac{1}{2}\,\nu\, r^2 \,(\FrameBasis\times \bm{r})\cdot \bm{b}_{\Section}\Big]\dA \\
 &=- \frac{1}{\Jsv}\int_{\Section} \Big[(\FrameBasis\times \bm{r})\cdot (\nabla \WarpShearShape{}^T
+ \tfrac{1}{2}\,\nu\, r^2 \,\oneS) \bm{b}_{\Section}\Big]\dA \\
 &= \FrameBasis\cdot\frac{1}{\Jsv}\left[\int_{\Section} \bm{r}\times (\nabla \WarpShearShape{}^T
+ \tfrac{1}{2}\,\nu\, r^2 \,\oneS)\dA \right] \bm{b}_{\Section} \\
\end{aligned}
\]
\end{Details}



\subsection{General solution}

Combining \eqref{eq:iesan-displacement-1} and \eqref{eq:iesan-displacement-2} furnishes the particular solution for the complete Saint-Venant class:
\begin{equation}
\begin{aligned}
\hat{\bm u}_{\rm P}%\{\bm a,\bm{b}_{\Section}\}
&=\Big(-\tfrac12 \reflenx^2\,\mathbf{1}-\nu\,\operatorname{dev}\bm{r}\otimes\bm{r}+\reflenx\,\FrameBasis\otimes\bm{r}\Big)\bm{a}_{\Section}
+(\reflenx\,\FrameBasis-\nu\,\bm{r})\,a_{\varepsilon}
+\big(\reflenx\,\FrameBasis\times\bm{r}+\WarpTwistIesan\,\FrameBasis\big)a_\vartheta\\
&\quad-\tfrac16\,\reflenx^3\,\bm{b}_{\Section}
-\reflenx\,\nu\Big(\operatorname{dev}\bm{r}\otimes\bm{r}-\bm{r}\otimes\CentroidLocation\Big)\bm{b}_{\Section}
+\tfrac12 \reflenx^2\,\FrameBasis\otimes(\bm{r}-\CentroidLocation)\,\bm{b}_{\Section}
+\FrameBasis\otimes\bm{\WarpShearIesan}(\bm{r})\,\bm{b}_{\Section}\\
&\quad+\big(\reflenx\,\FrameBasis\times\bm{r}+\WarpTwistIesan\,\FrameBasis\big)c_\vartheta, 
\end{aligned}
\label{eq:iesan-general-solution}
\end{equation}
%
\begin{Details}
\begin{equation*}
\begin{aligned}
\hat{\bm u}_{\rm P}%\{\bm a,\bm{b}_{\Section}\}
&=
 \reflenx\,\FrameBasis\,a_{\varepsilon} \\
&\quad-\tfrac12\,\reflenx^2\,\bm{a}_{\Section} \\
&\quad-\tfrac16\,\reflenx^3\,\bm{b}_{\Section} \\
&\quad+\reflenx\,\FrameBasis\otimes\bm{r}\,\bm{a}_{\Section} \\
&\quad+\tfrac12 \reflenx^2\,\FrameBasis\otimes\bm{r}\,\bm{b}_{\Section} \\
&\quad+\reflenx\,\FrameBasis\times\bm{r}\,(a_\vartheta+c_{\vartheta}) \\
&\quad-\tfrac12 \reflenx^2\,\FrameBasis\otimes\CentroidLocation\,\bm{b}_{\Section} \\
% \hline
&\quad+\bm{W}_{(\varepsilon)}\,a_{\varepsilon}
+ \bm{W}_{(a)}\,\bm{a}_{\Section} \\
&\quad+\WarpTwistIesan\,\FrameBasis\,(a_\vartheta+c_{\vartheta}) \\
&\quad+\FrameBasis\otimes\bm{\WarpShearIesan}(\bm{r})\,\bm{b}_{\Section}
+\reflenx\,\bm{\Pi}_{(b)}\bm{b}_{\Section} \\
%
&=
 \reflenx\,\FrameBasis\,a_{\varepsilon}
-\tfrac12\,\reflenx^2\,\bm{a}_{\Section} 
-\tfrac16\,\reflenx^3\,\bm{b}_{\Section} \\
&\quad+\reflenx\,\FrameBasis\times\bm{r}\,(a_\vartheta+c_{\vartheta}) \\
&\quad+\reflenx\quad\bm{r}\times\,(\FrameBasis\times\bm{a}_{\Section}) \\
&\quad+\tfrac12 \reflenx^2\,\bm{r}\times\,(\FrameBasis\times\bm{b}_{\Section}) \\
&\quad-\tfrac12 \reflenx^2\,\FrameBasis\otimes\CentroidLocation\,\bm{b}_{\Section} \\
&\quad+\bm{w}(x,\bm{r})
\end{aligned}
\end{equation*}
\end{Details}
with \(\bm a\) fixed by \eqref{eq:2-15}, \(\bm b\) by \eqref{eq:2-18}, \(\WarpTwistIesan\) by \eqref{eq:bvp-warp-twist-iesan},
\(\bm{\WarpShearIesan}\) by \eqref{eq:bvp-warp-shear-iesan},  \(c_\vartheta\) by \eqref{eq:iesan-2-22}. 
% and $\hat{\bm u}_{\rm R}$ is an arbitrary infinitesimal rigid body displacement field. 
%
The complete Saint-Venant solution is the sum of the particular solution \eqref{eq:iesan-general-solution}, 
and a homogeneous solution \(\hat{\bm u}_{\rm R}\):
\begin{equation}
\hat{\bm{u}}(x,\bm{r}) = \hat{\bm{u}}_{\rm P}(x,\bm{r}) + \hat{\bm{u}}_{\rm R}(x,\bm{r})
\label{eq:general-solution}
\end{equation}
and \(\hat{\bm{u}}_{\rm R}\) is an arbitrary the rigid body motion with six free parameters \(\bm{q}_u\in \mathbb{R}^3\) and \(\bm{q}_\theta\in \mathbb{R}^3\):
\begin{equation}
\begin{aligned}
\hat{\bm u}_{\rm R}\{\bm{q}\} 
&= \bm{q}_u + \bm{q}_\theta \times \hat{\bm{x}}(\reflenx,\bm{r}) \\
% &= \bm{q}_u + \bm{q}_\theta \times (\reflenx \FrameBasis + \bm{r})
\end{aligned}
\label{eq:rigid-field}
\end{equation}
%
The infinitesimal strain tensor is
\[
\StrainTensor
= \PointAxialStrain\,\FrameBasis\otimes\FrameBasis
\;+\;\FrameBasis\stackrel{\rm s}{\otimes} \bm{\varepsilon}_{\gamma}
\;-\;\nu\,\PointAxialStrain\,\oneS,
\]
where recall \(\bm{a}\stackrel{\rm s}{\otimes}\bm{b} = \tfrac{1}{2} (\bm{a}\otimes\bm{b} + \bm{b}\otimes\bm{a})\) is the symmetrized bun product and:
\begin{subequations}
\begin{align}
\PointAxialStrain
&= a_{\varepsilon} + \bm{r}\cdot\bm{a}_{\Section}
+ x\,(\bm{r}-\CentroidLocation)\cdot\bm{b}_{\Section},
 \label{eq:iesan-point-axial-strain} \\
%
\PointShearStrain 
&= (a_{\vartheta} + c_{\vartheta})(\FrameBasis\times \bm{r} + \nabla\WarpTwistIesan) 
+\bm{\Pi}_{(b)}\bm{b}_{\Section}
+ (\nabla \bm{\WarpShearIesan})^{\top} \bm{b}_{\Section} 
\label{eq:iesan-point-shear-strain}
\end{align}
\label{eq:iesan-point-strain}
\end{subequations}

\begin{Details}
Decompose $\nabla \hat{\bm{u}} = \FrameBasis\otimes\hat{\bm{u}}' + \gradS\hat{\bm{u}}$.

\begin{equation}
\begin{aligned}
\hat{\bm{u}}'
&=\; \Big(-x\,\one + \FrameBasis\otimes\bm{r}\Big)\bm{a}_{\Section}
\;+\; \FrameBasis\,a_{\varepsilon}
\;+\; (\FrameBasis\times\bm{r})\,(a_\vartheta+c_\vartheta)
\\&\quad
-\tfrac12 x^2\,\bm{b}_{\Section}
\;-\;\nu\bigl(\dev(\bm{r}\otimes\bm{r}) - \bm{r}\otimes\CentroidLocation\bigr)\bm{b}_{\Section}
\;+\; x\,\FrameBasis\otimes(\bm{r}-\CentroidLocation)\bm{b}_{\Section} \\
\gradS\hat{\bm{u}}
&=\; \left(-\nu\,(\bm{r}\cdot\bm{a}_{\Section})\,\oneS \;+\; x\,\FrameBasis\otimes \bm{a}_{\Section}\right)
 - \nu\,a_{\varepsilon}\,\oneS
\\&\quad
+ x\,( \FrameBasis\times \oneS)\,a_\vartheta
\;+\; \FrameBasis\otimes(\nabla\WarpTwistIesan)\,a_\vartheta
\;+\; x\,( \FrameBasis\times \oneS)\,c_\vartheta
\\&\quad
-\; x\nu\,\bigl[(\bm{r}-\CentroidLocation)\cdot\bm{b}_{\Section}\bigr]\oneS
\;+\; \tfrac12 x^2\,\FrameBasis\otimes \bm{b}_{\Section}
\;+\; \FrameBasis\otimes\bigl(c_\vartheta\,\gradS\WarpTwistIesan\bigr)
\\&\quad
+\; \FrameBasis\otimes \nabla\big(\bm{\WarpShearIesan}\cdot \bm{b}_{\Section}\big).
\end{aligned}
\label{eq:grad-uhat-iesan}
\end{equation}


Combine these to write
\[
\begin{aligned}
\nabla \hat{\bm{u}}
&= \FrameBasis\otimes\FrameBasis\,\Big(a_{\varepsilon}+(\bm{r}\cdot\bm{a}_{\Section})
+ \reflenx \,\big((\bm{r}-\CentroidLocation)\cdot\bm{b}_{\Section}\big)\Big)
\\&\quad
+\FrameBasis\otimes\Big\{
-\reflenx\,\bm{a}_{\Section}
+ (\FrameBasis\times\bm{r})\,(a_\vartheta+c_\vartheta)
-\tfrac12 \reflenx^2\,\bm{b}_{\Section}
% \\&\qquad\qquad
-\nu\bigl(\dev(\bm{r}\otimes\bm{r})-\bm{r}\otimes\CentroidLocation\bigr)\bm{b}_{\Section}
\Big\}
\\&\quad
+ \FrameBasis\otimes\Big\{
~\reflenx \,\bm{a}_{\Section} + \tfrac12 \reflenx ^2\,\bm{b}_{\Section}
+ (\nabla\varphi)\,a_\vartheta
+ c_\vartheta\,\nabla\WarpTwistIesan
+ \nabla \big(\bm{\WarpShearIesan}\cdot\bm{b}_{\Section}\big)
\Big\}
\\&\quad
-\nu\Big[a_{\varepsilon}+(\bm{r}\cdot\bm{a}_{\Section})
+ \reflenx\, (\bm{r}-\CentroidLocation)\cdot\bm{b}_{\Section}\Big]\oneS
\\&\quad
+ \reflenx\,(\FrameBasis\times \oneS)\,(a_\vartheta+c_\vartheta),
\end{aligned}
\]
where the last line is planar-skew (pure infinitesimal rotation).

\[
\StrainTensor \;\equiv\; \sym\nabla\hat{\bm{u}}
\;=\; \tfrac12\big(\nabla\hat{\bm{u}}+(\nabla\hat{\bm{u}})^{\mathrm t}\big).
\]
Using that $\sym(\FrameBasis\times\oneS)=\mathbf{0}$  and collecting symmetric parts furnishes the infinitesimal strain tensor. 
\end{Details}
% \bm V_\Sigma
% &= (\FrameBasis\times\bm{r} + \gradS \WarpTwistIesan)\,(a_\vartheta+c_\vartheta)
% \;+\; \gradS \big(\bm{\WarpShearIesan}\cdot\bm{b}_{\Section}\big)
% \;-\;\nu\bigl(\dev(\bm{r}\otimes\bm{r})-\bm{r}\otimes\CentroidLocation\bigr)\bm{b}_{\Section}\\

% \[
% \begin{aligned}
% \StrainTensor
% &= \underbrace{\Big[a_{\varepsilon}+(\bm{r}\cdot\bm{a}_{\Section})
% + x\,\big((\bm{r}-\CentroidLocation)\cdot\bm{b}_{\Section}\big)\Big]}_{\displaystyle \varepsilon}
% \,(\FrameBasis\otimes\FrameBasis)
% \\[2pt]
% &\quad
% +\;\sym\Big\{\FrameBasis\otimes \bm V_\Sigma\Big\}
% \;-\;
% \nu\,\underbrace{\Big[a_{\varepsilon}+(\bm{r}\cdot\bm{a}_{\Section})
% + x\,\big((\bm{r}-\CentroidLocation)\cdot\bm{b}_{\Section}\big)\Big]}_{\displaystyle \varepsilon}
% \,\oneS ,
% \end{aligned}
% \]
% with the mixed (axial–transverse) vector
% \[
% \bm V_\Sigma
% = (\FrameBasis\times\bm{r})\,(a_\vartheta+c_\vartheta)
% \;-\;\nu\bigl(\dev(\bm{r}\otimes\bm{r})-\bm{r}\otimes\CentroidLocation\bigr)\bm{b}_{\Section}
% \;+\; (\nabla\WarpTwistIesan)\,a_\vartheta
% \;+\; c_\vartheta\,\gradS\WarpTwistIesan
% \;+\; \nabla\!\big(\bm{\WarpShearIesan}\cdot\bm{b}_{\Section}\big) .
% \]

% \subsubsection{Equilibrium}
For any superposition satisfying \(\StressTensor\,\bm n=\mathbf{0}\) on \(\BodyLateralSurface\) and \eqref{eq:saint-venant-stress},
the field equations reduce on each section to
\begin{subequations}
\label{eq:sv-equilibrium}
\begin{align}
\bm{\tau}'&=\mathbf{0} && \text{in }\Section,\\
\DivS\bm{\tau}+\AxialStress'&=0 && \text{in }\Section,\\
\bm{\tau}\cdot\bm\nu&=0 && \text{on }\SectionBoundary,
\end{align}
\end{subequations}
where the prime denotes derivative with respect to \(\reflenx\).
\subparagraph{S1}
The Saint-Venant equilibrium equations for stress resultants 
$\snm$ form a first-order initial value problem:
\begin{equation}
\left\{
\begin{aligned}
\snm'(x) & = \mathbb{J}\,\snm(x), \\
\snm(0) &= \bm{s}_0
\end{aligned}\right.
\qquad\text{with}\qquad 
\snm := (N,\ShearForce,T,\bm{M})
\label{eq:equilibrium-ode}
\end{equation}
where \(\mathbb{J}\) is the equilibrium generator:
\begin{equation}
\mathbb{J} \coloneq \begin{bmatrix}
 \mathbf{0} & \mathbf{0}\\
 -\FrameBasis^\times & \mathbf{0}
\end{bmatrix}
\label{eq:J-matrix}
\end{equation}
\begin{Details}
\[
\mathbb{J}^{\top}\mathbb{J} = \begin{bmatrix}
    \oneS & \mathbf{0} \\
    \mathbf{0} & \mathbf{0}
\end{bmatrix},
\quad\text{and}\quad 
\mathbb{J}\mathbb{J}^{\top} = \begin{bmatrix}
    \mathbf{0} & \mathbf{0} \\
    \mathbf{0} & \oneS
\end{bmatrix},
\]
\end{Details}
Integrating \eqref{eq:equilibrium-ode} in view of the nilpotency property \(\mathbb{J}^2 = \mathbf{0}\) furnishes:
\begin{equation*}
\snm(x) = (\mathbf{1} + \reflenx\mathbb{J})\,\snm(0).
\end{equation*}

The resultants \(\snm\) at \(x=0\) thus depend linearly on the parameters \((a_{\varepsilon},\, \bm{a}_{\Section},\, a_{\vartheta},\, \bm{b}_{\Section})\):
\begin{equation}
\snm(0) = \IesanStiff\IesanVector
\label{eq:snm0-in-p}
\end{equation}
where \(\IesanVector\in\mathbb{R}^6\) collects the six kinematic  parameters of the particular solution: 
\begin{equation}
\IesanVector \coloneq (a_{\varepsilon},\, \bm{a}_{\Section},\, a_{\vartheta},\, \bm{b}_{\Section})
\label{eq:def-iesan-vector}
\end{equation}
In view of \eqref{eq:2-15} and \eqref{eq:2-18}:
\begin{equation}
\IesanStiff = \begin{bmatrix}
EA & EA\CentroidLocation^\top & 0 & \mathbf{0} \\[4pt]
\mathbf{0} & \mathbf{0} & \mathbf{0} & E\IntRGoRG \\[4pt]
0 & \mathbf{0} & G\Jsv & \mathbf{0} \\[4pt]
-EA\FrameBasis\times\CentroidLocation & -\FrameBasis^\times E\IntRoR & \mathbf{0} & \mathbf{0}
\end{bmatrix}.
\label{eq:G0-neg-m}
\end{equation}
where the centroidal inertia tensor \(\IntRGoRG\) is defined: 
\begin{equation}
    \IntRGoRG = \IntRoR - A\CentroidLocation\otimes\CentroidLocation
\label{eq:def-intRGoRG}
\end{equation}
\begin{Details}
Thus
\begin{equation}
\IesanStiff^{-1} =
\begin{bmatrix}
\displaystyle \frac{1}{EA} + \CentroidLocation\cdot (E\IntRGoRG)^{-1}\,\CentroidLocation
& \mathbf{0}
& 0
& -\CentroidLocation^\top (E\IntRGoRG)^{-1}\,\FrameBasis^\times
\\[6pt]
-(E\IntRGoRG)^{-1}\,\CentroidLocation
& \mathbf{0}
& \mathbf{0}
& (E\IntRGoRG)^{-1}\,\FrameBasis^\times
\\[6pt]
0
& \mathbf{0}
& \displaystyle \frac{1}{G\Jsv}
& \mathbf{0}
\\[6pt]
\mathbf{0}
& (E\IntRGoRG)^{-1}
& \mathbf{0}
& \mathbf{0}
\end{bmatrix}
\label{eq:G0-inv-neg-m}
\end{equation}
\end{Details}
Further, introduce the evolution operator \(\IesanEvolveSP\):
\begin{equation}
\IesanEvolveSP
:=\IesanStiff^{-1}\mathbb{J}\IesanStiff 
=\begin{bmatrix}
\mathbf{0} & \oneS - \FrameBasis\otimes\CentroidLocation \\[4pt]
\mathbf{0} & \mathbf{0}
\end{bmatrix}
= \begin{bmatrix}
0 & \mathbf{0} & 0 & -\CentroidLocation^{\top} \\[4pt]
\mathbf{0} & \mathbf{0} & \mathbf{0} & \oneS \\[4pt]
0 & \mathbf{0} & 0 & \mathbf{0} \\[4pt]
\mathbf{0} & \mathbf{0} & \mathbf{0} & \mathbf{0} \\
\end{bmatrix}
\label{eq:A-matrix-neg}
\end{equation}

\begin{Details}
Using $\IesanEvolveSP^2=\mathbf 0$:
\begin{equation*}
\snm(x)=\IesanStiff\exp\left(x\IesanEvolveSP\right)\bm{p} 
=\IesanStiff(\mathbf 1+x\IesanEvolveSP)\bm p,
% \label{eq:snm-in-p}
\end{equation*}
\end{Details}
\begin{Details}
Note:
\begin{equation*}
\mathbb{J}\IesanStiff = 
\begin{pmatrix}
    \mathbf{0} & \mathbf{0} \\
    \mathbf{0} & -\FrameBasis\times(E\IntRGoRG)
\end{pmatrix}
\end{equation*}
\end{Details}


\subsection{Summary}

\begin{table}[H]
    \centering
    \caption{Symbols employed for Saint~Venant's solution.}
    \label{tab:iesan-constants}
    \begin{tabular}{clcl}
    \toprule
    Symbol & Units & & Description \\ \midrule
        \(a_{\varepsilon}\in\mathbb{R}\) & $\sim$  & & Axial strain parameter \\
        \(a_{\vartheta}\in\mathbb{R}\)   & \Length{-1} & \\
        \(\bm{a}_{\Section}\in\Section\) & \Length{-1} & & Bending parameters for \(\mathrm{P}_{\rm I}\) \\
        \LightRule
        \(c_{\vartheta}\in\mathbb{R}\)   & \Length{-1} & & Torsion coupling under flexure \\ 
        \(b_{\varepsilon}\in\mathbb{R}\) & \\ %\Length{-1} & & \\
        \(\bm{b}_{\Section}\in\Section\) & \Length{-2} & \\
        \LightRule
        \(\bm{q}_u\in\mathbb{R}^3\) & L & & Rigid translation at \(x=0\) \\
        \(\bm{q}_{\theta}\in\mathbb{R}^3\) & & & Rigid rotation at \(x=0\), or constant transverse displacement rate   \\
        \LightRule
        \(\WarpTwistIesan\)   & \Length{2}  & \eqref{eq:bvp-warp-twist-iesan} \\
        \(\bm{\WarpShearIesan}\) & \Length{3}   & \eqref{eq:bvp-warp-shear-iesan} & \\
        \LightRule 
        \(\CentroidLocation\) & L & \eqref{eq:def-centroid} & Location of the centroid. \\
        \(\IntRoR\) & \Length{4} & \eqref{eq:def-j-origin} & \\
        \(\Jsv\) & \Length{4} & \eqref{eq:def-j-origin} & Torsion constant \\
        \(\CenterShearLocation\) & L & \eqref{eq:def-shear-center} & Torsion-flexure center \\
        \(\IntRGoRG\) & \Length{4} & \eqref{eq:def-intRGoRG} & Centroidal plane inertia tensor\\
        \LightRule 
        \(\snm\) & --- & \eqref{eq:equilibrium-ode} \\
        \(\IesanVector\) & --- & \eqref{eq:def-iesan-vector}\\
        \(\IesanStiff\) & --- & \eqref{eq:G0-neg-m} \\
    \bottomrule
    \end{tabular}
\end{table}
